\scanpage{66}
\deleted{Ceci donne une démonstration plus sympatique et élémentaire de 11.9.6. Quant à 11.9.3. (ii), il varésulter d'un

\uline{Lemme} (ou cor\deleted{l}ollaire à la proposition précédente). Sous les hypothèses de la proposition, ne supposons plus néc que $A'$ soit projectif sur $A$, mais seulement ceci: pour tout quotient intègre $B=A/\underline{p}$ de $A$, posant $B'=B\otimes_A A'$, il existe un $h\quad B-(0)$ tel que $B_h$ soit \deleted{projectif}\inserted{libre} sur $A_h$. Alors la conclusion de la proposition reste encore valable.

Soit en effet $\handwritten{\phi}$ une section de $F'$ dont les images par les $u'$ sont nulles, prouvons $\handwritten{\phi}=0$, i.e. $\operatorname{supp}\handwritten{\phi}=\handwritten{\unsure{\varnothing|\phi}}$. Dans le cas contraire, il existerait un point maximal $x'$ dans $\operatorname{supp}\handwritten{\phi}$, correspondant à l'idéal premier $\underline{p}'$ de $A'$. Soit $\deleted{\underline{p}=\underline{p}'\quad A'}$\inserted{$x$} son image dans $X$, défini par l'idéal $\underline{p}$ trace de $\underline{p}'$ sur $A$, alors en vertu de l'hypothèse appliquée à $A/\underline{p}$, on voit que, quitte à remplacer $X$ par un voisinage ouvert affine de $x$, on peut supposer que $\unsure{\deleted{B}A|\deleted{A}B}'/\underline{p}\unsure{\deleted{B}A|\deleted{A}B}'$ est \deleted{projectif}\inserted{libre} sur $A/\underline{p}$. Notons que nous pouvions choisir $x'$ de façon que $x$ soit égalem\unsure{\deleted{r}e|\deleted{e}r}nt maximal dans $g(\operatorname{supp}\handwritten{\phi})$ (commencer par choisir $x$ maximal dans $g(\operatorname{supp}\handwritten{\phi})$, puis $x'$ maixmal dans $g^{-1}(x)$), ce qui signifie que localisant en $\underline{p}$, le support de $\handwritten{\phi}_{\underline{p}}$ est contenu dans $\operatorname{Spec}(A'_{\underline{p}}/\underline{p}A'_{\underline{p}})$, donc $\handwritten{\phi}_{\underline{p}}$ est annulé par une puissance de $\underline{p}\ A'_{\underline{p}}$, ou encore que l'on peut trouver $\unsure{\deleted{u}h|\deleted{b}h|\deleted{h}b}\quad A-\underline{p}$ tel que}
